
\subsubsection{3.1 NLI Model}

The NLI model used is \texttt{cross-encoder/nli-deberta-v3-large} \citep{He2021DeBERTaV3}, a 304M-parameter cross-encoder trained on MNLI. Cross-encoders permit full token-level interaction between premise and hypothesis, enabling detection of subtle factual substitutions that would be missed by independent-encoding bi-encoders. The model is used frozen (inference-only; no gradient updates or fine-tuning).

\subsubsection{3.2 Net-Contradiction Scoring}

The hallucination score for a (context, response) pair is:

\begin{verbatim}
score = P(contradiction) − P(entailment)
\end{verbatim}

This framing amplifies the signal for commission-type hallucinations, which simultaneously increase P(contradiction) and decrease P(entailment). Pairs that are semantically neutral produce scores near zero, reducing false positives relative to using P(contradiction) alone.

\subsubsection{3.3 Sentence-Level Aggregation}

For dialogue and QA, the response is tokenized into sentences. Each (context, sentence) pair is scored independently, and the maximum net-contradiction score is taken as the response-level score. This follows the multiple-instance learning convention \citep{Laban2022SummaC}: a response is treated as hallucinated if any sentence contradicts the grounding context. For summarization, the full source document serves as the premise and the full summary as the hypothesis, scored at response level.

\subsubsection{3.4 Task-Specific Context Configuration}

\textbf{Dialogue}: A window of the three most recent dialogue turns is used as the premise. The rationale is that hallucinations in dialogue typically contradict recent context rather than the full history.

\textbf{QA}: The full knowledge passage is used as the premise; responses are scored at the response level.

\textbf{Summarization}: The full source document is used as the premise; the full summary is scored as the hypothesis.

\subsubsection{3.5 Design Scope and Ablation Status}

The three design choices described above --- net-contradiction framing, sentence-level max aggregation, and last-3-turn context windowing --- constitute the configuration evaluated in this study. Comparative ablation experiments (h-m2 through h-m4) that would isolate the contribution of each design choice were not executed. The results reported below characterize this specific configuration but cannot attribute performance differences to individual choices versus alternatives. This is acknowledged as Limitation L3 in Section 6.

\subsubsection{3.6 Dataset}

HaluEval \citep{Li2023HaluEval}: 20,000 balanced pairs per task (60,000 total across dialogue, QA, and summarization); binary labels (1 = hallucinated, 0 = non-hallucinated); hallucination examples generated by GPT-3.5. The dataset was loaded from \texttt{pminervini/HaluEval} (HuggingFace). Each raw example contains a \texttt{right\textbackslash \_X} / \texttt{hallucinated\textbackslash \_X} column pair; pairs were interleaved to produce balanced 20,000-pair sets with label 0 (correct) and label 1 (hallucinated), verified via \texttt{verify\textbackslash \_label\textbackslash \_map()}.

\subsubsection{3.7 Evaluation}

Primary metric: AUROC. Statistical significance: fastDeLong test \citep{DeLong1988Comparing} against a 0.5 null, $\alpha$ = 0.05. Gate criterion: AUROC > 0.55 with DeLong p < 0.05 on $\geq$ 2/3 tasks. Supplementary metrics: Cohen's d (effect size), Wilcoxon rank-sum test (score ordering), KL divergence from uniform (distributional non-uniformity).
